\section{Discussion}
\label{sec:discussion}

The central claim of this paper is that asking a human should be treated as part of robot planning, not as an external recovery behavior. This matters because the usefulness of a query depends on the robot's current belief and on the task consequences of remaining uncertain. A query about a visually ambiguous tomato is valuable if the answer changes whether the robot harvests, discards, or ignores it. The same query may be unnecessary if the tomato is irrelevant to the current goal or if another action can resolve the uncertainty autonomously.

The frontier belief representation is intended to make this connection explicit. By constructing uncertainty around states reachable from the current knowledge state and selected action, the robot avoids asking about facts that do not affect near-term planning. The entropy-based confidence score determines when the robot lacks enough information to commit a state, and the expected-entropy query rule determines which symbolic fact would most reduce the ambiguity.

Several limitations should be discussed clearly. First, the framework assumes that human answers are reliable or at least more reliable than the robot's uncertain observation. Future work should model noisy human feedback. Second, the current query form is binary and predicate-level. More expressive questions may reduce the number of interactions but increase user burden. Third, the planner uses a local frontier belief rather than a full global posterior. This improves tractability but may miss dependencies outside the immediate action-relevant frontier. Finally, the real-robot evaluation must show repeated trials rather than isolated demonstrations for the method to be convincing as a robotics contribution.
